% Complete Section 6.1, original PDF 73-78 / printed 63-68, including all Table 6.1 cells.
% block: exterior-algebra-heading page:73 type:section
\subsection{{Exterior algebra}}\label{sec:exterior-algebra}

% block: exterior-algebra-introduction page:73 type:ordinary-paragraph
{The primary objects of study in exterior calculus are differential $k$-forms, which we take up in the next section. These are functions whose values are alternating multilinear forms, also called algebraic $k$-forms. The exterior calculus is naturally concerned with the operations of differentiation and integration of these functions. But first we establish the basic terminology and properties of the algebraic $k$-forms which arise at each point, just as, before studying vector calculus, we must first understand finite dimensional vector spaces. This is the subject of exterior algebra, initiated by Grassman in the middle of the 19th century. Here, we briefly review the definitions and main properties, all of which are elementary.}

% block: multilinear-forms-heading page:73 type:subsection
\subsubsection{{Multilinear forms}}\label{subsec:multilinear-forms}

% block: multilinear-forms-definition pages:73,74 type:ordinary-paragraph
{For a finite dimensional vector space $V$, and a nonnegative integer $k$, we consider the vector space $\operatorname{Lin}^kV$ consisting of all $k$-linear maps}
\[
\omega:\overbrace{V\times\cdots\times V}^{k\text{ times}}\to\mathbb R.
\]
{By convention, we understand $\operatorname{Lin}^0V$ to be the space of scalars $\mathbb R$. If $\omega\in\operatorname{Lin}^jV$ and $\mu\in\operatorname{Lin}^kV$, the tensor product $\omega\otimes\mu\in\operatorname{Lin}^{j+k}V$ is given by}
\[
\begin{gathered}
(\omega\otimes\mu)(v_1,\ldots,v_{j+k})=\omega(v_1,\ldots,v_j)\mu(v_{j+1},\ldots,v_{j+k}),\\
v_1,\ldots,v_{j+k}\in V.
\end{gathered}
\]

% block: multilinear-dual-basis page:74 type:ordinary-paragraph
{From the definition, the space $\operatorname{Lin}^1V$ is just the dual space $V'$. We refer to its elements as either \emph{covectors} or $1$-forms. If we choose a basis $u_1,\ldots,u_n$ of $V$ ($n=\dim V$), this induces a dual basis $u^1,\ldots,u^n$ of $V'$, defined by $u^i(u_j)=\delta_{ij}$. A general element of $\operatorname{Lin}^kV$ is determined by the value it assigns to each of the $k$-tuples $(u_{\sigma_1},\ldots,u_{\sigma_k})$ with the $\sigma_i$ integers between $1$ and $n$. Thus $\dim\operatorname{Lin}^kV=n^k$. The element of $\operatorname{Lin}^kV$ which takes the $k$-tuple $(u_{\sigma_1},\ldots,u_{\sigma_k})$ to $1$ and all the other $k$-tuples $(u_{\tau_1},\ldots,u_{\tau_k})$ to $0$ is simply $u^{\sigma_1}\otimes\cdots\otimes u^{\sigma_k}$, and so the set of all of these gives a basis for $\operatorname{Lin}^kV$.}

% block: multilinear-canonical-coordinate-basis page:74 type:ordinary-paragraph
{The case $V=\mathbb R^n$ differs from the general case only in that there is a canonical basis $e_1,\ldots,e_n$. As notation for the corresponding dual basis we use $dx^1,\ldots,dx^n$. Thus, $dx^i$ is simply the linear functional which assigns to each vector $\mathbb R^n$ its $i$th component. The reason for the odd notation ($dx^i$ rather than, say, $e^i$) will become clear as we develop exterior calculus.}

% block: alternating-forms-heading page:74 type:subsection
\subsubsection{{Alternating multilinear forms}}\label{subsec:alternating-multilinear-forms}

% block: alternating-forms-definition page:74 type:ordinary-paragraph
{Multilinear forms $\omega\in\operatorname{Lin}^kV$ that change sign under the interchange of two variables are called \emph{skew-symmetric} or \emph{alternating}:}
\[
\begin{gathered}
\omega(v_1,\ldots,v_i,\ldots,v_j,\ldots,v_k)=-\omega(v_1,\ldots,v_j,\ldots,v_i,\ldots,v_k),\\
1\leq i<j\leq n,\qquad v_1,\ldots,v_k\in V.
\end{gathered}
\]
{In particular, an alternating form vanishes if any of its arguments is repeated. To any multilinear form $\omega$ there is associated the alternating form which is its skew-symmetric part:}
\[
(\operatorname{skw}\omega)(v_1,\ldots,v_k)=\frac1{k!}\sum_{\sigma\in\Sigma_k}\operatorname{sign}(\sigma)\omega(v_{\sigma_1},\ldots,v_{\sigma_k}),
\]
{where the sum is over all the permutations of the integers $1$ to $k$. For example, if $\omega\in\operatorname{Lin}^3V$, then}
\[
\begin{aligned}
(\operatorname{skw}\omega)(v_1,v_2,v_3)=\frac16[&\omega(v_1,v_2,v_3)+\omega(v_2,v_3,v_1)+\omega(v_3,v_1,v_2)\\
&-\omega(v_1,v_3,v_2)-\omega(v_3,v_2,v_1)-\omega(v_2,v_1,v_3)].
\end{aligned}
\]

% block: alternating-space-dimension page:74 type:ordinary-paragraph
{We define the space $\operatorname{Alt}^kV$ as the space of all skew-symmetric $k$-linear forms. The spaces $\operatorname{Alt}^0V$ and $\operatorname{Alt}^1V$ coincide with $\operatorname{Lin}^0V$ and $\operatorname{Lin}^1V$, i.e., are just the space of scalars and the space of covectors, respectively. An element of $\operatorname{Alt}^kV$ is determined by the value it assigns to each of the $k$-tuples $(e_{\sigma_1},\ldots,e_{\sigma_k})$ with $1\leq\sigma_1<\cdots<\sigma_k\leq n$, so $\dim\operatorname{Alt}^kV=\binom nk$. For $k>n$, $\operatorname{Alt}^kV=0$.}

% block: exterior-product-definition pages:74,75 type:ordinary-paragraph
{The \emph{exterior product} or \emph{wedge product} of $\omega\in\operatorname{Alt}^jV$ and $\mu\in\operatorname{Alt}^kV$ is given by}
\[
(\omega\wedge\mu)(v_1,\ldots,v_{j+k})=\sum_\sigma\operatorname{sign}(\sigma)\omega(v_{\sigma_1},\ldots,v_{\sigma_j})\mu(v_{\sigma_{j+1}},\ldots,v_{\sigma_{j+k}}),
\]
{where the sum is over $\sigma\in\Sigma_{j+k}$ for which $\sigma_1<\cdots<\sigma_j$ and $\sigma_{j+1}<\cdots<\sigma_{j+k}$. Equivalently,}
\[
\omega\wedge\mu=\binom{j+k}j\operatorname{skw}(\omega\otimes\mu).
\]
{The exterior product satisfies the anticommutativity law}
\[
\omega\wedge\mu=(-)^{jk}\mu\wedge\omega,\qquad\omega\in\operatorname{Alt}^jV,\ \mu\in\operatorname{Alt}^kV.
\]

% block: alternating-basis page:75 type:ordinary-paragraph
{Let $\{u_i\}$ be a basis for $V$ and $\{u^i\}$ the corresponding dual basis. For given $1\leq\sigma_1<\cdots<\sigma_k\leq n$, the element of $\operatorname{Alt}^kV$ which takes the $k$-tuple $(u_{\sigma_1},\ldots,u_{\sigma_k})$ to $1$ and takes all the other $k$-tuples with increasing indices to $0$ is}
\begin{equation}\label{eq:alternating-form-basis}
u^{\sigma_1}\wedge\cdots\wedge u^{\sigma_k}.
\end{equation}
{The set of all of such form a basis for $\operatorname{Alt}^kV$, and so the general alternating $k$-form on $V$ is a linear combination of these. In particular, the general alternating $k$-form on $\mathbb R^n$ can be written}
\[
\sum_{1\leq\sigma_1<\cdots<\sigma_k\leq n}a_\sigma\,dx^{\sigma_1}\wedge\cdots\wedge dx^{\sigma_k}
\]
{for a unique choice of coefficients $a_\sigma\in\mathbb R$.}

% block: alternating-operations-heading page:75 type:subsection
\subsubsection{{Operations on alternating forms}}\label{subsec:alternating-form-operations}

% block: alternating-contraction page:75 type:ordinary-paragraph
\medskip\noindent\textbf{{Contraction}}\par
{Besides the vector space operations on $\operatorname{Alt}^kV$ and the exterior product, there are several other useful algebraic operations. Given $\omega\in\operatorname{Alt}^kV$ with $k\geq1$ and $v\in V$, the \emph{contraction} of $\omega$ and $v$, also known as the interior product, is the $(k-1)$-form}
\[
(\omega\lrcorner v)(v_1,\ldots,v_{k-1})=\omega(v,v_1,\ldots,v_{k-1}).
\]
{(If $k=0$, $\omega\lrcorner v:=0$.) Since $\omega$ is alternating, $(\omega\lrcorner v)\lrcorner v$ always vanishes. Moreover, it is easy to check that contraction with a vector is an antiderivation in the sense that}
\[
\begin{gathered}
(\omega\wedge\eta)\lrcorner v=(\omega\lrcorner v)\wedge\eta+(-1)^k\omega\wedge(\eta\lrcorner v),\\
\omega\in\operatorname{Alt}^kV,\ \eta\in\operatorname{Alt}^lV.
\end{gathered}
\]

% block: alternating-inner-product page:75 type:ordinary-paragraph
\medskip\noindent\textbf{{Inner product}}\par
{If the vector space $V$ is endowed with an \emph{inner product} $\langle\cdot,\cdot\rangle$, then this induces the Riesz isomorphism $V$ to $V'$ and so an inner product on $V'$. An \emph{inner product} on $\operatorname{Alt}^kV$ is then defined by}
\[
\langle u^1\wedge\cdots\wedge u^k,v^1\wedge\cdots\wedge v^k\rangle=\det\big[(\langle u^i,v^j\rangle)_{i,j=1}^n\big],
\]
{which extends to general elements of $\operatorname{Alt}^kV$ by bilinearity. If the $u_i$ form an \emph{orthonormal} basis for $V$, then the corresponding dual basis elements $u^i$ are orthonormal in $V'$, as are the resulting basis functions \eqref{eq:alternating-form-basis} of $\operatorname{Alt}^kV$.}

% block: volume-form page:76 type:ordinary-paragraph
\medskip\noindent\textbf{{Volume form}}\par
{When $k$ is equal to the dimension $n$ of $V$, the space $\operatorname{Alt}^nV$ is one-dimensional, and so a basis consists of a single element. If the space $V$ is endowed with both an inner product and an \emph{orientation}, then there is a canonical choice of basis element. The ordered bases of a finite dimensional vector space $V$ divide into two classes according to whether the change-of-basis matrix connecting two of them has a positive or negative determinant. We orient the vector space by designating the ordered bases of one class as the positively oriented ones. (In $\mathbb R^3$, it is conventional to take as the positively oriented bases those which satisfy the "right-hand rule.") If $\omega\in\operatorname{Alt}^nV$, then the value of $\omega(u_1,\ldots,u_n)$ is the same for any orthonormal positively oriented basis $u_1,\ldots,u_n$, and we determine a unique element of $\operatorname{Alt}^nV$ by requiring this value to be $1$. The resulting alternating $n$-form is called the \emph{volume form} on $V$, which we typically denote by $\operatorname{vol}$. It may be expressed in terms of the dual basis to any positively oriented orthonormal basis as $\operatorname{vol}=u^1\wedge\cdots\wedge u^n$. On $\mathbb R^n$, the canonical basis $e_1,\ldots,e_n$ is orthonormal and positively oriented, so}
\[
\operatorname{vol}=dx^1\wedge\cdots\wedge dx^n.
\]

% block: hodge-star-definition page:76 type:ordinary-paragraph
\medskip\noindent\textbf{Hodge {star}}\par
{We continue with an oriented inner product space $V$. Let $\omega\in\operatorname{Alt}^kV$ for some $0\leq k\leq n$, and let $\mu\in\operatorname{Alt}^{n-k}V$. Then $\omega\wedge\mu=c\operatorname{vol}$ for some real number $c$, and the map $\mu\mapsto c$ is a linear functional on $\operatorname{Alt}^{n-k}$. By the Riesz representation theorem it is given by the inner product of $\mu$ with some other element of $\operatorname{Alt}^{n-k}V$. That element depends on $\omega$ and is called the \emph{Hodge star} or Hodge dual of $\omega$, denoted $\star\omega$. The Hodge star is thus characterized by the relation}
\begin{equation}\label{eq:hodge-star-characterization}
\omega\wedge\mu=\langle\star\omega,\mu\rangle\operatorname{vol},\qquad\omega\in\operatorname{Alt}^kV,\ \mu\in\operatorname{Alt}^{n-k}V.
\end{equation}

% block: hodge-star-basis-and-square page:76 type:ordinary-paragraph
{For each $k$, we thus have a linear map $\star:\operatorname{Alt}^kV\to\operatorname{Alt}^{n-k}V$. If we choose a positively oriented orthonormal basis of $1$-forms, $u^1,\ldots,u^n$, so $\operatorname{vol}=u^1\wedge\cdots\wedge u^n$, we find that}
\[
\star u^{\sigma_1}\wedge\cdots\wedge u^{\sigma_k}=\pm u^{\tau_1}\wedge\cdots\wedge u^{\tau_{n-k}}
\]
{where $\sigma$ and $\tau$ are increasing sequences which partition $\{1,\ldots,n\}$ and the sign is that of the permutation $(\sigma,\tau)$. For example, in $\mathbb R^3$,}
\[
\star dx^1=dx^2\wedge dx^3,\qquad\star dx^2=-dx^1\wedge dx^3,\qquad\star dx^3=dx^1\wedge dx^2.
\]
{Note that $\star$ is an \emph{isometry} and $\star\star\omega=(-)^{k(n-k)}\omega$, or, equivalently, if $\omega\in\operatorname{Alt}^kV$, then}
\[
\star\star\omega=\omega\text{ in odd dimensions},\qquad\star\star\omega=(-)^k\omega\text{ in even dimensions}.
\]
{Combining this with \eqref{eq:hodge-star-characterization} gives an alternate characterization of the Hodge star:}
\begin{equation}\label{eq:hodge-star-inner-product}
\omega\wedge\star\nu=\langle\omega,\nu\rangle\operatorname{vol},\qquad\omega,\nu\in\operatorname{Alt}^kV.
\end{equation}

% block: scalar-vector-proxies pages:76,77 type:ordinary-paragraph
\medskip\noindent\textbf{{Scalar and vector proxies}}\par
{If $V$ is an oriented inner product space of dimension $n$, then there is a natural isomorphism of each of the spaces $\operatorname{Alt}^0V$ and $\operatorname{Alt}^nV$ onto the scalar space $\mathbb R$, and of each of the spaces $\operatorname{Alt}^1V$ and $\operatorname{Alt}^{n-1}V$ onto the vector space $V$. In this way we can associate either a \emph{scalar proxy} or a \emph{vector proxy} to each member of $\operatorname{Alt}^kV$ for $k=0$, $1$, $n-1$, or $n$ (and so for all cases if $\dim V\leq3$). We now describe these isomorphisms.}

% block: proxy-isomorphisms page:77 type:ordinary-paragraph
{For any vector space $V$, $\operatorname{Alt}^0V$ is just the space $\mathbb R$ of scalars and the isomorphism is the identity. The inner product on $V$ defines the isomorphism of $\operatorname{Alt}^1V=V'$ onto $V$ via the Riesz representation. The Hodge star then maps $\operatorname{Alt}^{n-1}V$ isomorphically to $\operatorname{Alt}^1V$ (which is isomorphic to $V$) and similarly maps $\operatorname{Alt}^nV$ to $\operatorname{Alt}^0V=\mathbb R$.}

% block: two-dimensional-proxies page:77 type:ordinary-paragraph
{The case $n=2$ is particular, since in that case we have two ways of associating a vector proxy with an alternating $1$-form: either as a $1$-form by applying the Riesz map or as an $(n-1)$-form, that is, by applying first the Hodge star and then the Riesz map. These do not coincide, since the Hodge star $\operatorname{Alt}^1V\to\operatorname{Alt}^1V$ is not the identity map: it takes $u^1$ to $u^2$ and $u^2$ to $-u^1$. Therefore the two vector proxies differ by a rotation through a right angle.}

% block: linear-pullback page:77 type:ordinary-paragraph
\medskip\noindent\textbf{{Pullback under linear maps}}\par
{If $L:V\to W$ is a linear map between finite dimensional vector spaces, then there is an obvious \emph{pullback} operation $L^*:\operatorname{Alt}^kW\to\operatorname{Alt}^kV$ given by}
\begin{equation}\label{eq:linear-pullback-definition}
L^*\omega(v_1,\ldots,v_k)=\omega(Lv_1,\ldots,Lv_k).
\end{equation}
{The pullback acts contravariantly: if $K:W\to X$ is another map, then $(K\circ L)^*=L^*\circ K^*$. It also respects the algebraic structure, including the exterior product:}
\begin{equation}\label{eq:linear-pullback-wedge}
L^*(\omega\wedge\mu)=L^*\omega\wedge L^*\mu.
\end{equation}

% block: pullback-inclusion-and-normal page:77 type:ordinary-paragraph
{In the case when $V$ is a subspace of $W$ and $L$ is the inclusion, the pullback of a $k$-form on $W$ is its restriction to $V\times\cdots\times V$. Assuming that $W$ has an inner product, we may pass to vector proxies to view the pullback $\operatorname{Alt}^1W\to\operatorname{Alt}^1V$ as a map from $W\to V$, which is nothing other than the orthogonal projection. Another example, is when $W$ is an oriented inner product space of dimension $n$ and $V$ an oriented subspace of dimension $n-1$. In that case, we may use proxies to view the pullback $\operatorname{Alt}^{n-1}W\to\operatorname{Alt}^{n-1}V$ as a map from $W\to\mathbb R$. Let us calculate this map. We start with a positive orthonormal basis $u_2,\ldots,u_n$ of $V$ and prepend to it the unit normal $u_1$ of $V$ in $W$ with the sign chosen so that $u_1,\ldots,u_n$ is positively oriented. Let $u^1,\ldots,u^n$ denote the dual basis of $W'$. The pullback of $u^1$ to $V'$ vanishes, while the pullbacks $v^2,\ldots,v^n$ of the remaining $u^i$ give the basis of $V'$ dual to $u_2,\ldots,u_n$. Since pullback respects the exterior product, it follows that}
\[
\begin{gathered}
L^*(u^2\wedge\cdots\wedge u^n)=v^2\wedge\cdots\wedge v^n=\operatorname{vol}_V,\\
L^*(u^1\wedge u^2\wedge\cdots\wedge\widehat{u^i}\wedge\cdots\wedge u^n)=0.
\end{gathered}
\]
{Transferred to proxies, the pullback map $L^*$ thus becomes}
\[
u_1\mapsto1,\qquad u_i\mapsto0,\qquad i=2,\ldots,n.
\]
{This map is simply the inner product of a vector in $W$ with the oriented unit normal $u_1$. To summarize, in terms of proxies,}
\Needspace{6\baselineskip}
\begin{itemize}
\item{the pullback of a $1$-form to a subspace is the orthogonal projection, and}
\item{the pullback of an $n-1$ form to a subspace of codimension $1$ is the inner product with the normal.}
\end{itemize}

% block: exterior-algebra-r3-table page:78 type:table
\begin{table}[htbp]\centering\footnotesize
\caption{{Exterior algebra of $\mathbb R^3$ in terms of scalar and vector proxy fields.}}\label{tab:exterior-algebra-r3}
\setlength{\tabcolsep}{4pt}
\begin{tabular}{l|l}\hline
\multicolumn{2}{c}{{Proxies}}\\\hline
$\operatorname{Alt}^0\mathbb R^3=\mathbb R$&$c\leftrightarrow c$\\
$\operatorname{Alt}^1\mathbb R^3\xrightarrow{\cong}\mathbb R^3$&$u_1dx^1+u_2dx^2+u_3dx^3\leftrightarrow u$\\
$\operatorname{Alt}^2\mathbb R^3\xrightarrow{\cong}\mathbb R^3$&$u_1dx^2\wedge dx^3-u_2dx^1\wedge dx^3+u_3dx^1\wedge dx^2\leftrightarrow u$\\
$\operatorname{Alt}^3\mathbb R^3\xrightarrow{\cong}\mathbb R$&$c\leftrightarrow c\,dx^1\wedge dx^2\wedge dx^3$\\\hline
\multicolumn{2}{c}{{Exterior product}}\\\hline
$\wedge:\operatorname{Alt}^1\mathbb R^3\times\operatorname{Alt}^1\mathbb R^3\to\operatorname{Alt}^2\mathbb R^3$&$\times:\mathbb R^3\times\mathbb R^3\to\mathbb R^3$\\
$\wedge:\operatorname{Alt}^1\mathbb R^3\times\operatorname{Alt}^2\mathbb R^3\to\operatorname{Alt}^3\mathbb R^3$&$\cdot:\mathbb R^3\times\mathbb R^3\to\mathbb R$\\\hline
\multicolumn{2}{c}{{Contraction with a vector $v\in\mathbb R^3$}}\\\hline
$\lrcorner v:\operatorname{Alt}^1\mathbb R^3\to\operatorname{Alt}^0\mathbb R^3$&$v\cdot:\mathbb R^3\to\mathbb R$\\
$\lrcorner v:\operatorname{Alt}^2\mathbb R^3\to\operatorname{Alt}^1\mathbb R^3$&$v\times:\mathbb R^3\to\mathbb R^3$\\
$\lrcorner v:\operatorname{Alt}^3\mathbb R^3\to\operatorname{Alt}^2\mathbb R^3$&$v:\mathbb R\to\mathbb R^3$\quad$(c\mapsto cv)$\\\hline
\multicolumn{2}{c}{{Inner product}}\\\hline
{Inner product on $\operatorname{Alt}^k\mathbb R^3$}&{Euclidean inner product on $\mathbb R$ or $\mathbb R^3$}\\\hline
\multicolumn{2}{c}{{Volume form}}\\\hline
$\operatorname{vol}=dx^1\wedge dx^2\wedge dx^3$&$(v_1,v_2,v_3)\mapsto\det[v_1|v_2|v_3]$\\\hline
\multicolumn{2}{c}{Hodge {star}}\\\hline
$\star:\operatorname{Alt}^0\mathbb R^3\to\operatorname{Alt}^3\mathbb R^3$&$\operatorname{id}:\mathbb R\to\mathbb R$\\
$\star:\operatorname{Alt}^1\mathbb R^3\to\operatorname{Alt}^2\mathbb R^3$&$\operatorname{id}:\mathbb R^3\to\mathbb R^3$\\
$\star:\operatorname{Alt}^2\mathbb R^3\to\operatorname{Alt}^1\mathbb R^3$&$\operatorname{id}:\mathbb R^3\to\mathbb R^3$\\
$\star:\operatorname{Alt}^3\mathbb R^3\to\operatorname{Alt}^0\mathbb R^3$&$\operatorname{id}:\mathbb R\to\mathbb R$\\\hline
\multicolumn{2}{c}{{Pullback by a linear map $L:\mathbb R^3\to\mathbb R^3$}}\\\hline
$L^*:\operatorname{Alt}^0\mathbb R^3\to\operatorname{Alt}^0\mathbb R^3$&$\operatorname{id}:\mathbb R\to\mathbb R$\\
$L^*:\operatorname{Alt}^1\mathbb R^3\to\operatorname{Alt}^1\mathbb R^3$&$L^{\mathrm T}:\mathbb R^3\to\mathbb R^3$\\
$L^*:\operatorname{Alt}^2\mathbb R^3\to\operatorname{Alt}^2\mathbb R^3$&$(\det L)L^{-1}:\mathbb R^3\to\mathbb R^3$\\
$L^*:\operatorname{Alt}^3\mathbb R^3\to\operatorname{Alt}^3\mathbb R^3$&$(\det L):\mathbb R\to\mathbb R$\quad$(c\mapsto c\det L)$\\\hline
\end{tabular}
\end{table}

% block: exterior-algebra-r3-summary page:78 type:ordinary-paragraph
\medskip\noindent\textbf{{Exterior algebra of $\mathbb R^3$}}\par
{Since the spaces $\operatorname{Alt}^0\mathbb R^n$ and $\operatorname{Alt}^n\mathbb R^n$ may be identified with $\mathbb R$ and the spaces $\operatorname{Alt}^1\mathbb R^n$ and $\operatorname{Alt}^{n-1}\mathbb R^n$ with $\mathbb R^n$, in the case of $n=3$ (or fewer) dimensions, all the operations of exterior algebra may be viewed as operations on scalars or vectors. These operations are tabulated in Table~\ref{tab:exterior-algebra-r3}, which is adapted from \BookCite{11}. We remark that the matrix $(\det L)L^{-1}$ which shows up as the pullback of a $2$-form in the penultimate line of the table is the adjugate matrix of $L$, which is defined (as the transpose of the cofactor matrix) even when $L$ is singular.}
